\cotefile{72}
\scanpage{82}

% L1–2, top
\[
\text{donc}\quad
\underbrace{\operatorname{lin}(h_{i+1},\ldots,h_n)}_{Z_i\ \dim n-i\unsure{-1|+1}}
 \subset Z\subset
\underbrace{H_0\cap\cdots\cap H_i}_{Z'_i\ \dim n-i}.
\]

% L3–5
\deleted{Celà implique donc $Z$} et que cette condition implique bien que $Z$ est stable \unsure{par|pour} les $\unsure{\tau_j|\mathrm T_j}$. Elle \deleted{\unsure{ex}}implique aussi

% L6
\[
\deleted{Z.}\qquad Z_i\subset Z'_i
\]

% L7–9
ce qui équivaut aussi à

% L8–9, centre
\[
\begin{aligned}
 h_{i+1}&\notin H_i\\
 \text{i.e.}\qquad \beta_i&=0
\end{aligned}
\]

% L6–10, right
\margin{Dégager, pour $i$ fixé conditions équiv.:\\
a) $Z_i\subset Z'_i$\\
b) $Z_i$ inv. par $G$\\
c) $Z'_i$ inv. par $G$\\
d) $\beta_i=0$}

% L10–11
Et par une raison de dim, on doit avoir

% L11
\[
Z=Z_i\ \text{ou}\ Z=Z'_i.
\]

% L12
L'inverse étant clair.

% L13–16
\uline{Remarque} Conditions équivalentes, lorsque $\dim C\geqslant 2$ i.e $n\geqslant 2$ : a) Il n'y a pas d'hyperplan de $C$ invariant par $G$ b) On \deleted{\unsure{a}} a \deleted{\unsure{pas}} $\beta_0\ne0$, et \deleted{\unsure{est}} \inserted{on n'a pas}

% L16–19, underbrace and interline
\[
\underbrace{\left(\beta_1\deleted{\ne}=0\
 \deleted{\text{et}}\ 
 \overset{\inserted{\text{(i.e. }H_0\cap H_1\ \text{inv}\deleted{\unsure{.}}\ \deleted{\unsure{\text{ou}}}\ \text{et)}}}
 {\underbrace{H_0\cap H_1}_{\dim n-1}\subset C}\right)}
 _{\substack{\Updownarrow\ \text{si }\beta_0\ \text{ou }\mu_0\ne0\\
 \beta_1=0\ \text{et }\mu_1\mu_0-\beta_0=0}}.
\]

% archive sheet number: 81
